\section{Tests et validation}

\subsection{Stratégie de test}

Les tests doivent valider séparément chaque brique avant de valider la chaîne
complète. Cette approche réduit le risque de chercher une panne dans tout le
système alors qu'elle vient d'un seul module.

\subsection{Plan de tests}

\begin{longtable}{|p{3cm}|p{4cm}|p{4cm}|p{2cm}|}
\hline
\textbf{Test} & \textbf{Procédure} & \textbf{Résultat attendu} & \textbf{Statut} \\
\hline
Capture audio & Lancer une détection et vérifier la création du WAV temporaire. &
Fichier lisible généré par \texttt{arecord}. & À compléter \\
\hline
Reconnaissance AudD & Envoyer un extrait connu à l'API. &
Titre et artiste retournés. & À compléter \\
\hline
État sans résultat & Tester un extrait silencieux ou non reconnu. &
Affichage \textit{No result}. & À compléter \\
\hline
Affichage setup & Démarrer sans Wi-Fi connecté. &
QR code de configuration affiché. & À compléter \\
\hline
Affichage résultat & Simuler ou obtenir une réponse AudD. &
Titre, artiste, album, année et QR code affichés. & À compléter \\
\hline
Interface web & Ouvrir l'adresse IP du prototype sur le port \texttt{5000}. &
Page de configuration visible. & À compléter \\
\hline
Scan Wi-Fi & Appeler \texttt{/scan}. &
Liste JSON de réseaux disponibles. & À compléter \\
\hline
Bouton mode & Appuyer sur le bouton GPIO 5. &
Passage entre setup et running. & À compléter \\
\hline
Bouton action & Appuyer sur le bouton GPIO 6 en mode running. &
Une écoute de 10 secondes est lancée. Chaque nouvelle écoute demande un nouvel appui. & À compléter \\
\hline
Démo complète & Bouton, capture, reconnaissance, affichage. &
Chaîne fonctionnelle démontrable. & À compléter \\
\hline
\end{longtable}

\subsection{Feuille de recette}

\begin{longtable}{|p{8cm}|p{3cm}|p{3cm}|}
\hline
\textbf{Critère} & \textbf{Validé} & \textbf{Commentaire} \\
\hline
Le prototype démarre sans erreur bloquante. & Oui / Non & À compléter \\
\hline
Le mode initial correspond à l'état du Wi-Fi. & Oui / Non & À compléter \\
\hline
L'écran e-paper affiche les états principaux. & Oui / Non & À compléter \\
\hline
La capture audio produit un fichier exploitable. & Oui / Non & À compléter \\
\hline
La reconnaissance retourne un résultat sur un morceau connu. & Oui / Non & À compléter \\
\hline
Le système gère l'absence de résultat. & Oui / Non & À compléter \\
\hline
L'interface Wi-Fi est accessible. & Oui / Non & À compléter \\
\hline
Les boutons déclenchent les actions attendues. & Oui / Non & À compléter \\
\hline
La documentation permet de comprendre le fonctionnement. & Oui / Non & À compléter \\
\hline
\end{longtable}

\subsection{Limites connues}

\begin{itemize}
    \item La reconnaissance dépend d'une connexion Internet et du service AudD.
    \item Le token API doit être sécurisé avant toute diffusion publique du code.
    \item Les performances exactes sur Raspberry Pi doivent être mesurées sur le montage final.
    \item Le portail Wi-Fi n'est pas encore décrit comme un vrai portail captif complet ; il s'agit d'une interface locale Flask.
    \item Le boîtier, les schémas et le budget doivent être complétés avec les fichiers réels du prototype.
\end{itemize}

\subsection{Améliorations possibles}

\begin{itemize}
    \item déplacer la configuration sensible dans un fichier \texttt{.env} ;
    \item ajouter des tests automatisés unitaires pour les contrôleurs ;
    \item prévoir un mode simulation sans matériel pour la soutenance ;
    \item améliorer la mise en page e-paper avec des polices mieux adaptées ;
    \item ajouter un historique des derniers morceaux reconnus ;
    \item améliorer le boîtier et l'intégration mécanique ;
    \item mesurer l'autonomie et optimiser les cycles de rafraîchissement.
\end{itemize}
